% test-diagram-styles.tex — drie diagram-stijlen naast elkaar.
%
% Identiek model "N(t) = R*t" met R als directe instroom (constante),
% gerenderd in:
%   - tight     (default)            : valve krijgt R-label; R verdwijnt als losse node
%   - forrester (System Dynamics-std) : valve heeft geen label; R blijft staan; causale pijl
%   - edu       (didactisch dubbel)   : valve krijgt R-label én R blijft staan; causale pijl

\documentclass{article}
\usepackage[syntax=english]{numodel}

\begin{document}

\section*{tight (default)}
\numodelsetup{diagram-style=tight}
\newmodelprefix{a}
\mvar{T}{t}{0}{\s}{2}{systeem}
\mvar{Dt}{dt}{1}{\s}{2}{systeem}
\mvar{N}{n}{0}{}{0}{voorraad}
\mvar{R}{r}{5}{\per\s}{2}{constante}
\mrule{N}{\aN + \aR * \aDt}
\mrule{T}{\aT + \aDt}
\mstop{\aT >= 5}
\graphicmodel
\computemodel

\section*{forrester}
\numodelsetup{diagram-style=forrester}
\newmodelprefix{b}
\mvar{T}{t}{0}{\s}{2}{systeem}
\mvar{Dt}{dt}{1}{\s}{2}{systeem}
\mvar{N}{n}{0}{}{0}{voorraad}
\mvar{R}{r}{5}{\per\s}{2}{constante}
\mrule{N}{\bN + \bR * \bDt}
\mrule{T}{\bT + \bDt}
\mstop{\bT >= 5}
\graphicmodel
\computemodel

\section*{edu}
\numodelsetup{diagram-style=edu}
\newmodelprefix{c}
\mvar{T}{t}{0}{\s}{2}{systeem}
\mvar{Dt}{dt}{1}{\s}{2}{systeem}
\mvar{N}{n}{0}{}{0}{voorraad}
\mvar{R}{r}{5}{\per\s}{2}{constante}
\mrule{N}{\cN + \cR * \cDt}
\mrule{T}{\cT + \cDt}
\mstop{\cT >= 5}
\graphicmodel
\computemodel

\end{document}
